\section{Metric estimates at the torus scale}
\label{sec:metric-estimates}

This section and the next construct a special Lagrangian torus near a
maximal crossing of a polarized semistable degeneration. The coordinate
isotopy class is essential: an angular coordinate torus is parameterized
consistently around the puncture, so restriction to it is invariant under
monodromy. Section~\ref{sec:rotation} uses this identity in degree two.
The construction applies to Calabi--Yau manifolds independently of the
hyperk\"ahler setting. We use $d$ for the complex dimension and $N=2d$
for the real dimension.

\begin{theorem}[A special Lagrangian torus near a maximal crossing]
\label{thm:torus}
Let $\pi:\mathcal Y\to\Delta$ be the restriction to a disk of a projective
algebraic family over a smooth curve.  Assume that $\mathcal Y$ is smooth,
$\mathcal Y_0$ is a reduced simple normal crossings divisor, and the fibers
$Y_t$ for $t\ne0$ are smooth connected projective manifolds of complex
dimension $d\ge2$ with trivial canonical bundle.  Let $a_t$ be polarizing
classes given by a line bundle on $\mathcal Y$; this line bundle need not
be relatively ample at $0$.  Let $\eta_t$ be nowhere vanishing holomorphic
$d$-forms which extend as a section of the relative logarithmic canonical
line bundle.  Put $s=-\log|t|$, and suppose that, for constants $c,C>0$,
\begin{equation}
 c s^d\le I_t:=\int_{Y_t}|\eta_t|^2\le C s^d,
 \qquad 0<|t|\ll1.
 \label{eq:log-volume-growth}
\end{equation}
Suppose, moreover, that at some point $o\in\mathcal Y_0$ exactly $d+1$
components meet and that the logarithmic coefficient of $\eta$ at $o$ is
nonzero.  Thus in a coordinate neighborhood of $o$,
\begin{equation}
 z_0z_1\cdots z_d=t,
 \qquad
 \eta_t=f(z)\,\frac{\dd z_1}{z_1}\wedge\cdots\wedge
                         \frac{\dd z_d}{z_d},
 \qquad f(o)\ne0.
 \label{eq:crossing-chart}
\end{equation}
Equip $Y_t$ with the Ricci-flat K\"ahler metric in $a_t$.
For every sufficiently small nonzero $t$, the fiber $Y_t$ contains a
smooth embedded special Lagrangian real $d$-torus isotopic in $Y_t$ to an
angular coordinate torus in \eqref{eq:crossing-chart}.
Here special Lagrangian means that the K\"ahler form restricts to zero and
that $\eta_t$ has constant phase on the oriented torus.
\end{theorem}

We begin with an arbitrary sequence of nonzero parameters tending to zero.
The compactness argument will produce a subsequence on which every
sufficiently late fiber contains the required torus.  At the end of
Section~\ref{sec:completion}, a sequence of counterexamples will rule out
failure anywhere in a sufficiently small punctured disk.

This section chooses centers with noncollapsed unit balls and controls
the metric and logarithmic coordinates on every ball about them.
Section~\ref{sec:completion} identifies the pointed limit as a smooth
flat cylinder and perturbs its angular torus.  All constants below are
independent of $t$ in a sufficiently small punctured disk, unless stated
otherwise.  Every extraction is from the arbitrary initial sequence.

\subsection{Coordinate volume and mixed intersections}

Let $g_t^{(0)}$ and $\omega_t^{(0)}$ be the Ricci-flat metric and K\"ahler
form supplied by the Calabi--Yau theorem \cite{Yau}, with
$[\omega_t^{(0)}]=a_t$.  The useful scale is
\begin{equation}
 g_t=s g_t^{(0)},\qquad \omega_t=s\omega_t^{(0)},\qquad
 \mu_t=\dd V_t=\frac{\omega_t^d}{d!},\qquad
 \mu_t(Y_t)=c_a s^d,
 \quad c_a=\frac1{d!}\int_{Y_t}a_t^d>0.
 \label{eq:torus-scaling}
\end{equation}
The last constant is independent of $t$ by topological local triviality
over the punctured disk.  Scaling does not change either of the special
Lagrangian conditions.

Take $z_1,\ldots,z_d$ to be algebraic local equations of the corresponding
components.  They extend to rational functions on the algebraic total
space.  The remaining unit in the local equation of $\pi$ can be absorbed
into $z_0$; no rationality of this last coordinate will be used.  Shrink
the analytic coordinate neighborhood so that $|z_i|<r_*<1$ there.  On its
smooth fibers set
\[
 x_i=\log|z_i|,\qquad y_i=\arg z_i\in\R/(2\pi\Z),\qquad
 \gamma=\sum_{i=1}^d(\dd x_i^2+\dd y_i^2),\qquad
 \omega_\gamma=\sum_{i=1}^d\dd x_i\wedge\dd y_i.
\]
Thus $x_i+iy_i$ are local holomorphic logarithms, and $\gamma$ is the
standard flat metric on $\R^d\times(\R/2\pi\Z)^d$.  Its angular
circles have fixed length; we will compare $g_t$ with this metric.
Let
\[
 \Sigma=\left\{b\in\R^d:b_i<0\ (1\le i\le d),\quad
                              \sum_i b_i>-1\right\}.
\]
For a relatively compact subset $B\Subset\Sigma$, write
\[
 \mathcal T_t(B)=\{x/s\in B,\quad y\in(\R/2\pi\Z)^d\}.
\]
For small $t$ this tube is contained in the coordinate neighborhood.
Indeed all $|z_i|$, including
$|z_0|=\exp(-s-\sum_i x_i)$, are bounded by $e^{-\delta_B s}$
for some $\delta_B>0$.  In particular, these tubes converge uniformly to
$o$ in the total space.

Use the convention
$|\eta_t|^2=2^{-d}i^{d^2}\eta_t\wedge\overline{\eta_t}$, so that
\eqref{eq:crossing-chart} gives
$|\eta_t|^2=|f(z)|^2\dd V_\gamma$ in the chart.  Ricci-flatness implies
that its volume form is a constant multiple of $|\eta_t|^2$; integration
determines that constant.  Hence
\begin{equation}
 \dd V_t=\frac{c_a s^d}{I_t}|\eta_t|^2,
 \qquad
 \frac{\dd V_t}{\dd V_\gamma}
       =\frac{c_a s^d}{I_t}|f(z)|^2
       \longrightarrow\vartheta>0
 \quad\hbox{uniformly on }\mathcal T_t(B)
 \label{eq:coordinate-density}
\end{equation}
after a subsequence.  To obtain the convergence, first extract a limit
of $c_a s^d/I_t$ using \eqref{eq:log-volume-growth}, and then use
$f(z)\to f(o)\ne0$.  In particular, the two volume measures are uniformly
comparable on each fixed interior tube.

\begin{lemma}[Mixed-intersection estimates]
\label{lem:mixed-estimates}
For every fixed $B\Subset\Sigma$ and $0\le k\le d$,
\begin{equation}
 \int_{\mathcal T_t(B)}\omega_\gamma^k\wedge\omega_t^{d-k}
       \le C_B s^d.
 \label{eq:mixed-estimates}
\end{equation}
Consequently, with traces taken for real Riemannian metrics,
\begin{equation}
 \int_{\mathcal T_t(B)}\tr_\gamma g_t\,\dd V_\gamma\le C_Bs^d,
 \qquad
 \int_{\mathcal T_t(B)}\tr_{g_t}\gamma\,\dd V_t\le C_Bs^d.
 \label{eq:trace-estimates}
\end{equation}
\end{lemma}

\begin{proof}
Resolve the finitely many rational maps $z_i:\mathcal Y\dashrightarrow
\mathbb P^1$ simultaneously.  We obtain a projective modification
$\varpi:\widehat{\mathcal Y}\to\mathcal Y$ and morphisms
$f_i:\widehat{\mathcal Y}\to\mathbb P^1$, with $\varpi$ an isomorphism
over their common regular locus, in particular over the coordinate
neighborhood being used.  After shrinking the punctured disk, its
nearby fibers $\widehat Y_t$ are smooth.  This follows from generic
smoothness for the fixed algebraic model; no resolution depending on
$s$ is required.

Fix a Fubini--Study form $\omega_{\rm FS}$ on $\mathbb P^1$.  For
$A_b([u:v])=[e^bu:v]$, define the smooth closed semipositive form
\[
 \alpha_{s,t}=\frac1s\int_0^s\sum_{i=1}^d
                f_i^* A_b^*\omega_{\rm FS}\,\dd b
 \quad\hbox{on }\widehat Y_t.
\]
Its cohomology class is the restriction of
$\sum_i f_i^*[\omega_{\rm FS}]$, independently of $s$, since each $A_b$
is homotopic to the identity.  Thus every intersection
\begin{equation}
 \int_{\widehat Y_t}\alpha_{s,t}^k\wedge
                  (\varpi^*\omega_t^{(0)})^{d-k}
 \label{eq:fixed-mixed-intersection}
\end{equation}
is a fixed finite number.  Both classes involved extend over the fixed
model, so this number is independent of $t$ as well.

In a logarithmic coordinate $w_i=x_i+iy_i$, the $i$th summand before
averaging is a fixed positive constant times
\[
 \frac{e^{2(b+x_i)}}{(1+e^{2(b+x_i)})^2}
                   \dd x_i\wedge\dd y_i.
\]
For $x/s\in B$, both $-x_i$ and $s+x_i$ are at least
$\delta_Bs$.  Consequently the interval
$[-x_i-1,-x_i+1]$ lies in $[0,s]$ for large $s$.  On this interval the
displayed scalar factor is bounded below by a positive constant.
Summing over $i$ proves, as Hermitian forms on the tube,
\[
 \alpha_{s,t}\ge\frac{c_B}{s}\omega_\gamma.
\]
The pullback $\varpi^*\omega_t^{(0)}$ is semipositive, including on the
exceptional divisors.  Positivity of wedge products therefore gives
\[
 \begin{split}
 \int_{\mathcal T_t(B)}\omega_\gamma^k\wedge\omega_t^{d-k}
 &\le C_Bs^k s^{d-k}
       \int_{\widehat Y_t}\alpha_{s,t}^k\wedge
                        (\varpi^*\omega_t^{(0)})^{d-k}\\
 &\le C_B' s^d.
 \end{split}
\]
This proves \eqref{eq:mixed-estimates}.  Finally, for Hermitian metrics
$h,k$ with K\"ahler forms $\omega_h,\omega_k$, the real trace identity is
\[
 (\tr_h k)\,\dd V_h
       =\frac{2}{(d-1)!}\omega_k\wedge\omega_h^{d-1}.
\]
Use $k=1$ and $k=d-1$ in \eqref{eq:mixed-estimates} to obtain the two
assertions in \eqref{eq:trace-estimates}.
\end{proof}

\subsection{Global energies and many noncollapsed centers}

The two trace bounds have different roles.  The bound for
$\tr_{g_t}\gamma$ controls logarithmic functions on intrinsic balls;
the bound for $\tr_\gamma g_t$ supplies many noncollapsed centers.
We also record two energies whose averages tend to zero.

Choose concentric open Euclidean balls
$B_0\Subset B_1\Subset B_2\Subset\Sigma$, and a smooth function
$0\le\rho\le1$ supported compactly in $B_2$ and identically one on $B_1$.
We will select centers in $\mathcal T_t(B_0)$.  Define functions on
the entire compact fiber by extending the following by zero:
\begin{equation}
 U_0=s\rho(x/s),\qquad U_j=x_j\rho(x/s)\quad(1\le j\le d),
 \qquad H_t=\sum_{j=0}^d|\nabla U_j|_{g_t}^2.
 \label{eq:cutoff-functions}
\end{equation}
These extensions are smooth because their support is compactly
contained in the chart for each fixed $t$.  On this support $x_j/s$ is
bounded, and
\[
 \dd U_0=\sum_i\rho_i(x/s)\,\dd x_i,
 \qquad
 \dd U_j=\rho(x/s)\,\dd x_j
           +\frac{x_j}{s}\sum_i\rho_i(x/s)\,\dd x_i.
\]
It follows that $H_t\le C\tr_{g_t}\gamma$ on the support and that
\begin{equation}
 \frac1{\mu_t(Y_t)}\int_{Y_t}H_t\,\dd V_t\le C_H.
 \label{eq:cutoff-global-energy}
\end{equation}

Choose a fixed projective embedding of the family over a neighborhood
of the disk, written
$(\Phi,\pi):\mathcal Y\hookrightarrow\mathbb P^M\times\Delta$.
Give $\mathbb P^M$ a fixed Fubini--Study metric and set
\[
 K_t=|\Rm_{g_t}|^2+|\dd\Phi_t|_{g_t}^2.
\]
The Ricci-flat K\"ahler Chern--Weil identity expresses
$\int_{Y_t}|\Rm_{g_t^{(0)}}|^2\dd V_{g_t^{(0)}}$ as a
universal dimension-dependent positive constant times
$\int_{Y_t}c_2(Y_t)\,a_t^{d-2}$.  This is independent of $t$.
The energy of the holomorphic map $\Phi_t$ in the unscaled metric is
likewise a fixed multiple of
$\int_{Y_t}\Phi_t^*[\omega_{\rm FS}]\,a_t^{d-1}$.
Under $g_t=s g_t^{(0)}$, squared curvature scales by $s^{-2}$,
map-energy density by $s^{-1}$, and volume by $s^d$.  Therefore
\begin{equation}
 \begin{split}
 \int_{Y_t}|\Rm_{g_t}|^2\,\dd V_t&\le C s^{d-2},\\
 \int_{Y_t}|\dd\Phi_t|_{g_t}^2\,\dd V_t&\le C s^{d-1},\\
 m_t:=\frac1{\mu_t(Y_t)}\int_{Y_t}K_t\,\dd V_t
        &\le C(s^{-2}+s^{-1})\longrightarrow0.
 \end{split}
 \label{eq:small-global-energies}
\end{equation}
This also covers $d=2$, when the curvature integral in the unscaled
metric is purely topological.

\begin{lemma}[A positive fraction of noncollapsed centers]
\label{lem:many-noncollapsed-centers}
There exist constants $\delta,v>0$ and measurable sets
$E_t\subset\mathcal T_t(B_0)$ such that
\begin{equation}
 \mu_t(E_t)\ge\delta\mu_t(Y_t),\qquad
 \mu_t(B_{g_t}(q,1))\ge v\quad(q\in E_t).
 \label{eq:many-noncollapsed-centers}
\end{equation}
All balls here and below use the intrinsic distance of the complete
manifold $(Y_t,g_t)$.
\end{lemma}

\begin{proof}
Fix a rectangular angle chart of positive size in $(\R/2\pi\Z)^d$.
The region $sB_0$ contains at least $c_0s^d$ disjoint translates of a
fixed small $d$-dimensional cube.  Taking their products with one fixed
angle cube gives at least $c_0s^d$ disjoint convex coordinate boxes $Q$,
all of the same Euclidean size, inside $\mathcal T_t(B_0)$.
By the first estimate of \eqref{eq:trace-estimates}, at least half of
these boxes satisfy
\begin{equation}
 \int_Q\tr_\gamma g_t\,\dd V_\gamma\le C_0
 \label{eq:good-box-trace}
\end{equation}
for a constant $C_0$ independent of $t$.  By
\eqref{eq:coordinate-density}, every such box has
$m_0\le\mu_t(Q)\le m_1$ for fixed positive $m_0,m_1$.

For $a,b\in Q$, the straight coordinate segment lies in $Q$, and its
length bounds the intrinsic distance from above.  Consequently
\[
 d_{g_t}(a,b)\le |a-b|_\gamma
   \int_0^1\sqrt{\tr_\gamma g_t((1-r)a+rb)}\,\dd r.
\]
Write $h=\sqrt{\tr_\gamma g_t}$.  For $0\le r\le1/2$, first integrate
in $a$ and change variables $z=(1-r)a+rb$.  Its Jacobian is
$(1-r)^N\ge2^{-N}$ and $z$ stays in $Q$.  Thus
\[
 \int_Q\int_Q h((1-r)a+rb)\,\dd a\,\dd b
       \le2^N\vol_\gamma(Q)\int_Qh\,\dd V_\gamma.
\]
For $1/2\le r\le1$, use $b$ first and the Jacobian $r^N$ instead.
The Euclidean diameter and volume of $Q$ are fixed.  Cauchy--Schwarz and
\eqref{eq:good-box-trace} now bound the double integral of
$d_{g_t}(a,b)$ uniformly.  Comparing measures gives a uniform constant
$C_1$ such that
\[
 \frac1{\mu_t(Q)^2}\int_Q\int_Qd_{g_t}(a,b)\,
                       \dd V_t(a)\,\dd V_t(b)\le C_1.
\]

The set $E_Q$ of $a\in Q$ for which
$\mu_t(Q)^{-1}\int_Qd_{g_t}(a,b)\dd V_t(b)\le2C_1$
has measure at least $\mu_t(Q)/2$.  For each such $a$, a second application
of Markov's inequality gives
\[
 \mu_t\bigl(B_{g_t}(a,R_0)\bigr)\ge\mu_t(Q)/2\ge m_0/2,
 \qquad R_0=\max\{1,4C_1\}.
\]
Since $\Ric(g_t)=0$, Bishop--Gromov comparison yields
$\mu_t(B_{g_t}(a,1))\ge R_0^{-N}m_0/2$.
Finally take the union of $E_Q$ over the disjoint good boxes.  There are
at least $(c_0/2)s^d$ of them and each contributes measure at least
$m_0/2$.  In view of \eqref{eq:torus-scaling}, this union has a fixed
positive fraction of the total volume, as required.
\end{proof}

\subsection{One center with estimates at every radius}

We now choose one of the noncollapsed centers so that the average cutoff
energy stays bounded, and the average curvature and embedding energies
tend to zero, on every ball about it.  The maximal function makes this
simultaneous choice possible.

For a nonnegative integrable function $h$ on $Y_t$, let
\[
 \mathcal M_t h(q)=\sup_{r>0}
       \frac1{\mu_t(B_{g_t}(q,r))}
                   \int_{B_{g_t}(q,r)}h\,\dd V_t
\]
be the centered maximal function.  Bishop--Gromov gives doubling with
constant $2^N$.  The usual disjoint-ball covering argument therefore
has a constant $D_N$ independent of $t$ and gives
\begin{equation}
 \mu_t\{\mathcal M_t h>\lambda\}
       \le\frac{D_N}{\lambda}\int_{Y_t}h\,\dd V_t.
 \label{eq:maximal-weak-type}
\end{equation}
For completeness, take at each point of the set on the left a ball
whose average exceeds $\lambda$, select disjoint balls whose fivefold
enlargements cover a compact subset of that set, and use
$\mu_t(5B)\le5^N\mu_t(B)$.  Sum the integrals on the disjoint balls
and exhaust the set by compact subsets.  Radii larger than the diameter
simply give the whole manifold, so no diameter bound is needed.

Choose $A$ so large that $D_NC_H/A<\delta/3$.  Set
$\epsilon_t=\sqrt{m_t}+s^{-1}$, which is positive and tends to zero.
The estimate \eqref{eq:maximal-weak-type}, first for $H_t$ and then for
$K_t$, shows that
\[
 \mu_t\{\mathcal M_tH_t>A\}<\frac\delta3\mu_t(Y_t),\qquad
 \mu_t\{\mathcal M_tK_t>\epsilon_t\}
       \le D_N\sqrt{m_t}\,\mu_t(Y_t).
\]
For small $t$ the latter quantity is also less than
$\delta\mu_t(Y_t)/3$.  These two exceptional sets cannot cover $E_t$.
We can therefore select $p_t\in E_t$ with
\begin{equation}
 \begin{gathered}
 \mu_t(B_{g_t}(p_t,1))\ge v,\\
 \frac1{\mu_t(B_{g_t}(p_t,R))}
       \int_{B_{g_t}(p_t,R)}H_t\,\dd V_t\le A,\\
 \frac1{\mu_t(B_{g_t}(p_t,R))}
       \int_{B_{g_t}(p_t,R)}K_t\,\dd V_t\le\epsilon_t
       \quad\hbox{for every }R>0.
 \end{gathered}
 \label{eq:good-center}
\end{equation}
After a further subsequence,
\[
 b_t:=x(p_t)/s\longrightarrow b_*\in\overline{B_0}\subset B_1.
\]
In particular $\rho(b_t)=1$,
$U_0(p_t)=s$, and $U_j(p_t)=x_j(p_t)$.

We next turn average gradient control into an estimate centered at the
actual value of the function at $p_t$.  The ball Poincar\'e inequality
for nonnegative Ricci curvature, a consequence of the segment inequality
\cite{CheegerColdingSegment}, states that
\[
 \frac1{\mu_t(B(q,r))}\int_{B(q,r)}|U-U_{B(q,r)}|^2\,\dd V_t
 \le C_Nr^2\frac1{\mu_t(B(q,2r))}
                       \int_{B(q,2r)}|\nabla U|^2\,\dd V_t.
\]
Here $U_{B(q,r)}$ denotes the mean and doubling absorbs any equivalent
choice of enlarged radius.  Applied at $p_t$ to $U_j$, this gives an
upper bound $CA r^2$.  Writing $a_r=(U_j)_{B(p_t,r)}$ and comparing the
smaller ball with the larger one by doubling, we also obtain
\[
 \begin{split}
 |a_{r/2}-a_r|^2
 &\le\frac1{\mu_t(B(p_t,r/2))}
           \int_{B(p_t,r/2)}|U_j-a_r|^2\,\dd V_t\\
 &\le C A r^2.
 \end{split}
\]
Since $U_j$ is smooth, $a_{2^{-k}R}\to U_j(p_t)$ as $k\to\infty$.
Summing the preceding bounds over dyadic radii yields
$|a_R-U_j(p_t)|\le C\sqrt A\,R$.  Combining this with Poincar\'e proves
\begin{equation}
 \frac1{\mu_t(B_{g_t}(p_t,R))}
   \int_{B_{g_t}(p_t,R)}|U_j-U_j(p_t)|^2\,\dd V_t
       \le C R^2
 \quad(0\le j\le d,\ R>0).
 \label{eq:centered-poincare}
\end{equation}

The cutoff functions now locate each fixed intrinsic ball in the
crossing chart.  Divide \eqref{eq:centered-poincare} by $s^2$ and use
$b_t\to b_*$.  On every fixed-radius ball, $U_0/s\to1$ and
$U_j/s\to(b_*)_j$ in $L^2$.  Where $U_0/s>1/2$, the point lies in the
coordinate tube and
\[
 \frac{x_j}{s}=\frac{U_j/s}{U_0/s}.
\]
Chebyshev's inequality consequently gives, for every $R,\eta>0$,
\begin{equation}
 \mu_t\left(B_{g_t}(p_t,R)\setminus
     \{q\hbox{ in the chart}:|x(q)/s-b_*|<\eta\}\right)
          \longrightarrow0.
 \label{eq:coordinate-concentration}
\end{equation}
The assertion concerns absolute volume, because Bishop--Gromov gives
the uniform bound
$\mu_t(B_{g_t}(p_t,R))\le\omega_NR^N$, where $\omega_N$ is the volume
of the Euclidean unit ball.

\subsection{The regular part of the pointed limit}

Noncollapse permits us to pass to a limit.  Vanishing curvature energy
will make its regular set flat; vanishing embedding energy will place
each regular compact set inside the chosen crossing chart.

The noncollapse in \eqref{eq:good-center} and $\Ric(g_t)=0$ give, after
extraction, a pointed measured Gromov--Hausdorff limit
\[
 (Y_t,g_t,p_t,\mu_t)\longrightarrow(Z,d_Z,p,\mu).
\]
We use the following precise consequences of noncollapsed Ricci limit
theory and Einstein regularity \cite{CheegerColding,ColdingVolume,CheegerNaber}.
The space $Z$ is complete and proper; $\mu$ is its $N$-dimensional
Riemannian volume on its open dense regular set $G$.  Regular points
are the points with a Euclidean tangent cone.  Compact subsets of $G$
admit smooth convergence identifications with regions in $Y_t$, under
which the metrics converge in $C^\infty$.  The closed singular set
$S=Z\setminus G$ satisfies
\begin{equation}
 \dim_{\rm H} S\le N-4,
 \qquad \mu(B_Z(q,r))\le\omega_Nr^N.
 \label{eq:limit-singular-set}
\end{equation}
The codimension assertion is Theorem~1.4 of \cite{CheegerNaber}; the
smooth convergence follows from Einstein $\epsilon$-regularity
\cite[Theorem~2.11]{CheegerNaber} and elliptic bootstrapping.
Complex structures also converge smoothly on
$G$: they are parallel orthogonal tensors for the converging metrics.
Their limit $I_\infty$ is parallel and integrable, so the limiting
metric $g$ is K\"ahler on $G$.

Every compact smooth region lies in the image of a fixed-radius ball
about $p_t$ for all sufficiently small $t$ along the sequence.  By
\eqref{eq:good-center} and the volume upper bound, its curvature energy
is at most $\epsilon_t\omega_NR^N$, which tends to zero.  Smooth metric
convergence therefore implies
\begin{equation}
 \Rm_g=0\quad\hbox{on }G.
 \label{eq:flat-regular-set}
\end{equation}
It remains to put these smooth regions in the coordinate chart about
$o$.  We use the following small-energy argument for the maps $\Phi_t$.

\begin{lemma}[Vanishing energy on a smooth domain]
\label{lem:small-energy}
Let $h_\nu$ be metrics converging smoothly on a fixed smooth domain $D$
of real dimension $N$,
and let $\psi_\nu:(D,h_\nu)\to P$ be smooth harmonic maps to a fixed
compact Riemannian manifold.  If their energies tend to zero on every
relatively compact subdomain, then
$\sup_K|\dd\psi_\nu|\to0$ for every $K\Subset D$.
After extraction they converge smoothly to a constant on each connected
relatively compact domain.
\end{lemma}

\begin{proof}
Smooth harmonic maps are stationary, which is the property used in the
energy monotonicity argument below.  Fix nested relatively
compact domains around $K$.  Their metrics have uniform smooth bounds,
and geodesic balls with centers in an intermediate domain have a
uniformly positive admissible radius $r_0$.
For a smooth harmonic map the stress tensor
$|\dd\psi|^2h-2\psi^*g_P$ has zero divergence.  Testing this identity
against radial vector fields and using
$\nabla^2(r^2/2)=h+O(r^2)$ in these geodesic balls gives
\[
 \Theta_\nu'(q,r)\ge-Cr\Theta_\nu(q,r),\qquad
 \Theta_\nu(q,r)=r^{2-N}
       \int_{B_{h_\nu}(q,r)}|\dd\psi_\nu|^2\,\dd V_{h_\nu}.
\]
The nonnegative radial derivative term has been discarded here.
Integration gives the uniform monotonicity consequence
\begin{equation}
 \Theta_\nu(q,r)\le C\Theta_\nu(q,r_0)
     \le Cr_0^{2-N}E_\nu\longrightarrow0
       \quad(0<r\le r_0),
 \label{eq:harmonic-map-monotonicity}
\end{equation}
where $E_\nu$ is the energy on a fixed larger domain.  This is the
stationary harmonic-map monotonicity argument; see also \cite{Lin}.

If derivatives were unbounded on $K$, maximize
$\operatorname{dist}_{h_\nu}(q,\partial D')
 |\dd\psi_\nu(q)|$ on a fixed intermediate domain $D'$.
At maximizing points $q_\nu$, put
$L_\nu=|\dd\psi_\nu(q_\nu)|$.  Then $L_\nu\to\infty$ and the product
just maximized tends to infinity.  On the ball of radius
$L_\nu^{-1}$ about $q_\nu$, point selection gives
$|\dd\psi_\nu|\le2L_\nu$ for large $\nu$.  Rescale the domain metric by
$L_\nu^2$.  The rescaled maps $v_\nu$ on unit balls satisfy
$|\dd v_\nu(0)|=1$ and $|\dd v_\nu|\le2$, while the rescaled domain
metrics converge smoothly to a Euclidean metric.  The harmonic-map
Bochner inequality and these derivative bounds imply
\[
 \Delta|\dd v_\nu|^2\ge-C|\dd v_\nu|^2.
\]
The local mean-value inequality, with uniform constants on these
smooth balls, now gives
\[
 1\le C\int_{B(0,1)}|\dd v_\nu|^2
   =C L_\nu^{N-2}
          \int_{B_{h_\nu}(q_\nu,L_\nu^{-1})}
                         |\dd\psi_\nu|^2\,\dd V_{h_\nu}
   =C\Theta_\nu(q_\nu,L_\nu^{-1})\longrightarrow0,
\]
a contradiction.  Thus the derivatives are uniformly bounded on
interior domains.  Apply the same Bochner and mean-value argument on
fixed smaller balls: it gives
$\sup_K|\dd\psi_\nu|^2\le C E_\nu\to0$.
Compactness of $P$ and interior elliptic estimates for the smooth
harmonic-map equation then give smooth subsequential convergence; the
limit has zero derivative and is constant on connected domains.
\end{proof}

The maps $\Phi_t$ are holomorphic between K\"ahler manifolds and hence
harmonic.  On a compact subdomain of $G$, pull them back using the
smooth convergence identifications.  Their energies tend to zero by
\eqref{eq:good-center}, so Lemma~\ref{lem:small-energy} applies.
Every local constant limit is $\Phi(o)$.  Indeed a fixed smaller regular
ball has positive limiting volume, whereas the exceptional set in
\eqref{eq:coordinate-concentration} has volume tending to zero.
That ball therefore contains points with $x/s$ in a fixed compact
subset of $\Sigma$.  Such points converge to $o$ in the total space,
and their images converge to $\Phi(o)$.  Local uniform convergence of
the maps identifies the constant on each regular ball separately,
regardless of whether $G$ is connected.
It follows, by a finite cover, that
\begin{equation}
 \Phi_t\longrightarrow\Phi(o)
       \quad\hbox{uniformly on every compact subset of }G.
 \label{eq:embedding-map-limit}
\end{equation}

Since $(\Phi,\pi)$ is an embedding, there is a neighborhood of
$(\Phi(o),0)$ whose inverse image lies in the coordinate neighborhood
\eqref{eq:crossing-chart}.  Thus the corresponding regions of $Y_t$
are eventually in that chart.  On such a region all $|z_i|<r_*<1$, and
$\sum_{i=0}^d\log|z_i|=-s$.  For $1\le j\le d$ this gives
\[
 -s-d\log r_*\le x_j\le\log r_*.
\]
In particular $x_j/s$ is uniformly bounded there.  Each $x_j$ is
pluriharmonic, being locally the real part of $\log z_j$, and is
therefore harmonic for $g_t$.  Its boundedness and
\eqref{eq:coordinate-concentration} give local $L^2$ convergence of
$x_j/s$ to $(b_*)_j$ on the smooth regions.  Interior harmonic estimates
for the smoothly converging metrics improve this to
\begin{equation}
 \frac{x_j}{s}\longrightarrow(b_*)_j
       \quad\hbox{in }C^\infty_{\rm loc}(G),\qquad 1\le j\le d.
 \label{eq:regular-chart-convergence}
\end{equation}
As $b_*\in B_1$ and $\rho=1$ there, every compact subset of $G$ has a
neighborhood on whose corresponding regions $\rho(x/s)=1$ for all
sufficiently small $t$.  Hence on these regions $U_j=x_j$ for $j\ge1$.
The all-radius energy bound \eqref{eq:good-center} and the centered
oscillation estimate \eqref{eq:centered-poincare} therefore apply to
the actual logarithmic coordinates.  In the next section they give a
global coordinate map on the limit, including its possible singular
points.
